% GENERATED FILE. Edit canonical lessons, then run npm run generate:books.
% canonical-source-hash: fnv1a64:e5c9ae1cd0a60570
% canonical-lessons: MR-C516-sanganak, MR-C516-bhramandhvani, MR-C516-durdhvani, MR-C516-citraphit, MR-C516-vaparne

\chapter{Screens and Phones}
\label{ch:mr-screens-and-phones}
\hlchaptermodality{516}

\begin{chapteropening}
\textbf{By the end of this chapter:} \emph{I can talk about computers, phones and video clips.}

Complete the last lesson of chapter 516: I can talk about computers, phones and video clips.
\end{chapteropening}

\section[\texorpdfstring{\mr{संगणक} (saṅgaṇak)}{संगणक (saṅgaṇak)}]{\mr{संगणक} (saṅgaṇak) — a computer}

\label{lesson:MR-C516-sanganak}



\begin{quote}
Before the new one: say the Marathi for an accountant, then the Marathi for to wrap up.
\end{quote}

\subsection*{You'll want to know: \mr{संगणक}}
\textbf{\mr{संगणक}} — \emph{saṅgaṇak} — \textquotedblleft{}a computer\textquotedblright{}.

\begin{quote}
\mr{माझ्या} \mr{भावाने} \mr{नवीन} \mr{संगणक} \mr{घेतला}.
\end{quote}

\begin{quote}
\emph{mājhyā bhāvāne navīn saṅgaṇak ghetlā.}
\end{quote}

\begin{quote}
My brother bought a new computer.
\end{quote}

Masculine.

\begin{cousinweb}
\textbf{\mr{संगणक}} (\emph{saṅgaṇak}) is a modern Marathi coinage built from Sanskrit \emph{saṁ-}, \textquotedblleft{}together\textquotedblright{}, and \emph{gaṇaka}, \textquotedblleft{}one who counts\textquotedblright{} (Monier-Williams): the machine that reckons.
\end{cousinweb}

\begin{grammarlens}[title={What you've built}]
One more word for describing what you lived and felt.
\end{grammarlens}

\subsection*{Your turn}
\begin{itemize}
\raggedright
  \item \emph{Say it:} \emph{saṅgaṇak}
  \item \emph{Say it:} \emph{saṅgaṇak}, once more
  \item \emph{Say it:} \emph{mājhyā bhāvāne navīn saṅgaṇak ghetlā.} — My brother bought a new computer.
  \item \emph{Recall:} say the Marathi for an accountant, then the Marathi for to wrap up, then say \emph{saṅgaṇak} again
\end{itemize}

\begin{tcolorbox}[breakable,colback=teal!4,colframe=teal!35!black,title={Before you move on}]
What does \mr{संगणक} mean? (\textquotedblleft{}A computer\textquotedblright{}.) Say it once more.
\end{tcolorbox}
\section[\texorpdfstring{\mr{भ्रमणध्वनी} (bhramaṇdhvanī)}{भ्रमणध्वनी (bhramaṇdhvanī)}]{\mr{भ्रमणध्वनी} (bhramaṇdhvanī) — a mobile phone}

\label{lesson:MR-C516-bhramandhvani}



\begin{quote}
Before the new one: say the Marathi for to wrap up, then the Marathi for a computer.
\end{quote}

\subsection*{You'll want to know: \mr{भ्रमणध्वनी}}
\textbf{\mr{भ्रमणध्वनी}} — \emph{bhramaṇdhvanī} — \textquotedblleft{}a mobile phone\textquotedblright{}.

\begin{quote}
\mr{मीराचा} \mr{भ्रमणध्वनी} \mr{बसमध्ये} \mr{हरवला}.
\end{quote}

\begin{quote}
\emph{mīrācā bhramaṇdhvanī basmadhye haravlā.}
\end{quote}

\begin{quote}
Meera's mobile phone got lost on the bus.
\end{quote}

Masculine.

In speech people say \mr{मोबाईल}; \mr{भ्रमणध्वनी} is the formal word, a sound that roams.

\begin{cousinweb}
\textbf{\mr{भ्रमणध्वनी}} (\emph{bhramaṇdhvanī}) is a modern tatsama coinage from Sanskrit \emph{bhramaṇa}, \textquotedblleft{}roaming\textquotedblright{}, and \emph{dhvani}, \textquotedblleft{}sound\textquotedblright{} (Monier-Williams).
\end{cousinweb}

\begin{grammarlens}[title={What you've built}]
One more word for describing what you lived and felt.
\end{grammarlens}

\subsection*{Your turn}
\begin{itemize}
\raggedright
  \item \emph{Say it:} \emph{bhramaṇdhvanī}
  \item \emph{Say it:} \emph{bhramaṇdhvanī}, once more
  \item \emph{Say it:} \emph{mīrācā bhramaṇdhvanī basmadhye haravlā.} — Meera's mobile phone got lost on the bus.
  \item \emph{Recall:} say the Marathi for to wrap up, then the Marathi for a computer, then say \emph{bhramaṇdhvanī} again
\end{itemize}

\begin{tcolorbox}[breakable,colback=teal!4,colframe=teal!35!black,title={Before you move on}]
What does \mr{भ्रमणध्वनी} mean? (\textquotedblleft{}A mobile phone\textquotedblright{}.) Say it once more.
\end{tcolorbox}
\section[\texorpdfstring{\mr{दूरध्वनी} (dūrdhvanī)}{दूरध्वनी (dūrdhvanī)}]{\mr{दूरध्वनी} (dūrdhvanī) — a telephone; a landline}

\label{lesson:MR-C516-durdhvani}



\begin{quote}
Before the new one: say the Marathi for a computer, then the Marathi for a mobile phone.
\end{quote}

\subsection*{You'll want to know: \mr{दूरध्वनी}}
\textbf{\mr{दूरध्वनी}} — \emph{dūrdhvanī} — \textquotedblleft{}a telephone; a landline\textquotedblright{}.

\begin{quote}
\mr{कार्यालयातला} \mr{दूरध्वनी} \mr{सकाळपासून} \mr{वाजत} \mr{होता}.
\end{quote}

\begin{quote}
\emph{kāryālayātlā dūrdhvanī sakāḷpāsūn vājat hotā.}
\end{quote}

\begin{quote}
The office telephone had been ringing since morning.
\end{quote}

Masculine.

\begin{cousinweb}
\textbf{\mr{दूरध्वनी}} (\emph{dūrdhvanī}) is a modern tatsama coinage from Sanskrit \emph{dūra}, \textquotedblleft{}far\textquotedblright{}, and \emph{dhvani}, \textquotedblleft{}sound\textquotedblright{} (Monier-Williams).
\end{cousinweb}

\begin{grammarlens}[title={What you've built}]
One more word for describing what you lived and felt.
\end{grammarlens}

\subsection*{Your turn}
\begin{itemize}
\raggedright
  \item \emph{Say it:} \emph{dūrdhvanī}
  \item \emph{Say it:} \emph{dūrdhvanī}, once more
  \item \emph{Say it:} \emph{kāryālayātlā dūrdhvanī sakāḷpāsūn vājat hotā.} — The office telephone had been ringing since morning.
  \item \emph{Recall:} say the Marathi for a computer, then the Marathi for a mobile phone, then say \emph{dūrdhvanī} again
\end{itemize}

\begin{tcolorbox}[breakable,colback=teal!4,colframe=teal!35!black,title={Before you move on}]
What does \mr{दूरध्वनी} mean? (\textquotedblleft{}A telephone; a landline\textquotedblright{}.) Say it once more.
\end{tcolorbox}
\section[\texorpdfstring{\mr{चित्रफीत} (citraphīt)}{चित्रफीत (citraphīt)}]{\mr{चित्रफीत} (citraphīt) — a video clip}

\label{lesson:MR-C516-citraphit}



\begin{quote}
Before the new one: say the Marathi for a mobile phone, then the Marathi for a telephone.
\end{quote}

\subsection*{You'll want to know: \mr{चित्रफीत}}
\textbf{\mr{चित्रफीत}} — \emph{citraphīt} — \textquotedblleft{}a video clip\textquotedblright{}.

\begin{quote}
\mr{मी} \mr{सहलीची} \mr{एक} \mr{छोटी} \mr{चित्रफीत} \mr{बनवली}.
\end{quote}

\begin{quote}
\emph{mī sahlīcī ek choṭī citraphīt banavlī.}
\end{quote}

\begin{quote}
I made a short video of the trip.
\end{quote}

Feminine.

The oblique is \mr{चित्रफिती}-: \mr{चित्रफितीत}, in the video.

\begin{cousinweb}
\textbf{\mr{चित्रफीत}} (\emph{citraphīt}) is a Marathi compound of \mr{चित्र}, \textquotedblleft{}picture\textquotedblright{}, from Sanskrit \emph{citra} (Monier-Williams), and \mr{फीत}, \textquotedblleft{}ribbon, tape\textquotedblright{}, a loan from Portuguese \emph{fita}.
\end{cousinweb}

\begin{grammarlens}[title={What you've built}]
One more word for describing what you lived and felt.
\end{grammarlens}

\subsection*{Your turn}
\begin{itemize}
\raggedright
  \item \emph{Say it:} \emph{citraphīt}
  \item \emph{Say it:} \emph{citraphīt}, once more
  \item \emph{Say it:} \emph{mī sahlīcī ek choṭī citraphīt banavlī.} — I made a short video of the trip.
  \item \emph{Recall:} say the Marathi for a mobile phone, then the Marathi for a telephone, then say \emph{citraphīt} again
\end{itemize}

\begin{tcolorbox}[breakable,colback=teal!4,colframe=teal!35!black,title={Before you move on}]
What does \mr{चित्रफीत} mean? (\textquotedblleft{}A video clip\textquotedblright{}.) Say it once more.
\end{tcolorbox}
\section[\texorpdfstring{\mr{वापरणे} (vāparṇe)}{वापरणे (vāparṇe)}]{\mr{वापरणे} (vāparṇe) — to use}

\label{lesson:MR-C516-vaparne}



\begin{quote}
Before the new one: say the Marathi for a telephone, then the Marathi for a video clip.
\end{quote}

\subsection*{You'll want to know: \mr{वापरणे}}
\textbf{\mr{वापरणे}} — \emph{vāparṇe} — \textquotedblleft{}to use\textquotedblright{}.

\begin{quote}
\mr{आजोबांनी} \mr{पहिल्यांदा} \mr{भ्रमणध्वनी} \mr{वापरला}.
\end{quote}

\begin{quote}
\emph{ājobānnī pahilyāndā bhramaṇdhvanī vāparlā.}
\end{quote}

\begin{quote}
Grandfather used a mobile phone for the first time.
\end{quote}

Past, with -\mr{ने} (he …): \mr{त्याने} \mr{संगणक} \mr{वापरला}, \emph{tyāne saṅgaṇak vāparlā}.

\begin{cousinweb}
\textbf{\mr{वापरणे}} (\emph{vāparṇe}) is usually traced to Sanskrit \emph{vyāpārayati}, \textquotedblleft{}employs\textquotedblright{}. CDIAL lists under it the related Marathi \emph{vāvarṇẽ}, \textquotedblleft{}to manage, to do the business of\textquotedblright{} (CDIAL 12206), but not \mr{वापरणे} itself, so the path is not settled.
\end{cousinweb}

\begin{grammarlens}[title={What you've built}]
One more word for describing what you lived and felt.
\end{grammarlens}

\subsection*{Your turn}
\begin{itemize}
\raggedright
  \item \emph{Say it:} \emph{vāparṇe}
  \item \emph{Say it:} \emph{vāparṇe}, once more
  \item \emph{Say it:} \emph{ājobānnī pahilyāndā bhramaṇdhvanī vāparlā.} — Grandfather used a mobile phone for the first time.
  \item \emph{Recall:} say the Marathi for a telephone, then the Marathi for a video clip, then say \emph{vāparṇe} again
\end{itemize}

\begin{tcolorbox}[breakable,colback=teal!4,colframe=teal!35!black,title={Before you move on}]
What does \mr{वापरणे} mean? (\textquotedblleft{}To use\textquotedblright{}.) Say it once more.
\end{tcolorbox}
